\section{Task 3 --- Graph Embedding over Homogeneous Structures}
\label{p1:sec:main}

\subsection{Approach}
\label{p1:sec:approach}

All four notebook TODOs are implemented; the provided scaffolding (\texttt{load\_data},
\texttt{adjacency\_lookups}, \texttt{evaluate\_ranking}, \texttt{breakdown\_by\_source\_degree},
\texttt{run\_experiment}, \texttt{comparison\_table}, \dots) is used verbatim and
\texttt{RANDOM\_SEED}${}=0$.  Loading reproduces the published statistics exactly:
$7{,}624$ users, $19{,}464$ observed edges, $666$ isolated users, mean degree $5.106$,
max degree $153$, $734$ components, and $4{,}171\times 21$ validation/test queries.

\paragraph{\texttt{sample\_walks}.}
A strategy is nothing but a transition rule; the \emph{walk budget is identical for every
strategy} ($10$ walks of length $40$ from each of the $6{,}958$ non-isolated users, giving
$69{,}580$ walks in every run).  The $666$ isolated users return \texttt{[]}, as the
docstring permits, and therefore keep a zero embedding row.  Seven rules are compared:

\begin{center}\small
\begin{tabular}{@{}llp{6.6cm}@{}}
\toprule
strategy & transition rule $P(x \mid \cdot)$ & hypothesis being tested \\
\midrule
DeepWalk & uniform over $N(v)$ & baseline: unbiased 1st-order walk \\
node2vec & $\propto 1/p,\,1,\,1/q$ by distance to $t$ & $q<1$ outward (homophily), $q>1$ inward (structural equivalence) \\
RWR & uniform, teleport to \texttt{start} w.p.\ $r$ & a strongly \emph{personalised} context should sharpen pairwise proximity, at the cost of global community signal \\
deg-corrected & $\propto \deg(x)^{-\alpha}$ & uniform walks over-visit hubs; equalising visit counts should lift the low-degree source buckets \\
triadic & $\propto 1 + |N(v)\cap N(x)|$ & preferring triangle-closing steps keeps the walk inside a community, i.e.\ maximal homophily \\
\bottomrule
\end{tabular}
\end{center}

\noindent
RWR, deg-corrected and triadic are my own additions.  node2vec uses the standard
$O(1)$ rejection sampler rather than materialising alias tables, which keeps a
second-order walk at $1.3\,$s for the whole corpus.

\paragraph{\texttt{train\_embedding}.}
Skip-Gram with negative sampling (\texttt{gensim.Word2Vec}, \texttt{sg=1}) over the walk
corpus.  I set \texttt{sample=0}, i.e.\ \emph{no} frequent-token subsampling: subsampling
would itself down-weight hubs and so half-implement the deg-corrected strategy,
confounding that comparison.  \texttt{workers=1} makes every run bit-exactly reproducible
(gensim's multi-threaded training is not deterministic).

\paragraph{Downstream models --- held fixed across all strategies.}
As the notebook demands, \texttt{LinkPredictor} and \texttt{CountryClassifier} are the
same code for every row of the comparison table; the embedding is the only thing that
changes.  Both $\ell_2$-normalise $Z$ row-wise (zero rows stay zero), standardise, and fit
a logistic regression.  A pair $(u,v)$ becomes
\[
  \phi(u,v)=\bigl[\,\underbrace{z_u \odot z_v}_{\text{Hadamard}},\;
                   \underbrace{|z_u-z_v|}_{\text{abs.\ diff.}},\;
                   \cos(z_u,z_v),\; z_u^\top z_v,\; \|z_u-z_v\|_2 \,\bigr]
  \in \mathbb{R}^{2d+3},
\]
the node2vec paper's two best binary operators concatenated with three scalar similarities.
\texttt{score} returns the continuous \texttt{decision\_function}, so higher always means
"more likely an edge".

\begin{table}[H]
\centering\small
\caption*{Fixed hyper-parameters (identical in every run).}
\begin{tabular}{@{}ll@{\hspace{2em}}ll@{}}
\toprule
walks per node & $10$                  & embedding dim.\ $d$ & $128$ \\
walk length    & $40$                  & Skip-Gram window    & $5$ \\
corpus size    & $69{,}580$ walks      & negative samples    & $5$ \\
seed           & $0$ (also $1,2$)      & epochs              & $5$ \\
LP regulariser & $C\!\in\!\{0.1,1,10\}$, chosen by val.\ MRR & subsampling & off (\texttt{sample=0}) \\
CC regulariser & $C\!\in\!\{0.1,1,10,100\}$, chosen by val.\ micro-F1 & solver & lbfgs \\
\bottomrule
\end{tabular}
\end{table}

\subsection{Implementation}
\label{p1:sec:code}

Only the code I wrote is listed; the provided scaffolding is omitted.

\lstset{basicstyle=\ttfamily\scriptsize,breaklines=true,language=Python,
        keywordstyle=\color{blue!60!black},commentstyle=\color{green!35!black},
        stringstyle=\color{red!50!black},showstringspaces=false,columns=fullflexible}

\begin{lstlisting}[caption={TODO 1 --- \texttt{sample\_walks}},captionpos=b]
def sample_walks(self, G, start):
    """`num_walks` walks of length `walk_length`, all beginning at `start`."""
    adj, deg = adjacency_lookups(G)
    if deg[start] == 0:                    # isolated node -> no walks (allowed)
        return []
    rng, L, R, strat = self._rng, self.walk_length, self.num_walks, self.strategy

    if strat == "node2vec":                # 2nd-order, O(1) rejection sampling
        NBR = self._nbr_lists(G)           # cached sorted neighbour lists
        p, q = self.params["p"], self.params["q"]
        inv_p, inv_q, wmax = 1/p, 1/q, max(1/p, 1.0, 1/q)
        walks = []
        for _ in range(R):
            walk = [start, rng.choice(NBR[start])]      # 1st step unbiased
            while len(walk) < L:
                cur, prev = walk[-1], walk[-2]
                cn, prev_nbrs = NBR[cur], adj[prev]
                while True:
                    x = rng.choice(cn)
                    w = inv_p if x == prev else (1.0 if x in prev_nbrs else inv_q)
                    if rng.random() * wmax < w:
                        break
                walk.append(x)
            walks.append(walk)
        return walks

    if strat in ("degcorr", "triadic"):    # 1st-order, edge-weighted
        tables = self._weighted_tables(G)  # per node: (nbrs, cumulative weights)
        walks = []                         # degcorr: w(x) = deg(x)**-alpha
        for _ in range(R):                 # triadic: w(x) = 1 + |N(cur) & N(x)|
            walk = [start]
            while len(walk) < L:
                nbrs, cum = tables[walk[-1]]
                walk.append(rng.choices(nbrs, cum_weights=cum, k=1)[0])
            walks.append(walk)
        return walks

    NBR = self._nbr_lists(G)               # "uniform" (DeepWalk) and "rwr"
    restart = self.params.get("restart", 0.0) if strat == "rwr" else 0.0
    walks = []
    for _ in range(R):
        walk = [start]
        while len(walk) < L:
            if restart and rng.random() < restart:
                walk.append(start)         # teleport home, then carry on
            else:
                walk.append(rng.choice(NBR[walk[-1]]))
        walks.append(walk)
    return walks
\end{lstlisting}

\begin{lstlisting}[caption={TODO 2 --- \texttt{train\_embedding}},captionpos=b]
def train_embedding(self, walks, n_nodes):
    d = self.params.get("dim", 128)
    model = Word2Vec(sentences=[[str(v) for v in w] for w in walks],
                     vector_size=d, window=self.params.get("window", 5),
                     sg=1, hs=0, negative=self.params.get("negative", 5),
                     epochs=self.params.get("epochs", 5),
                     min_count=0, sample=0, seed=self.seed, workers=1)
    Z = np.zeros((n_nodes, d), dtype=float)   # nodes never walked keep a zero row
    Z[[int(k) for k in model.wv.index_to_key]] = model.wv.vectors
    return Z
\end{lstlisting}

\begin{lstlisting}[caption={TODO 3a --- \texttt{LinkPredictor} (and the pair feature map)},captionpos=b]
def _unit(Z):                       # row-wise L2 norm; all-zero rows stay zero
    n = np.linalg.norm(Z, axis=1, keepdims=True)
    return Z / np.where(n == 0.0, 1.0, n)

def _pair_features(Z, pairs):       # (n_q, n_cand, 2) -> (n_q*n_cand, 2d+3)
    u, v = Z[pairs[..., 0]].reshape(-1, Z.shape[1]), Z[pairs[..., 1]].reshape(-1, Z.shape[1])
    had, dif = u * v, np.abs(u - v)
    dot = had.sum(1, keepdims=True)
    l2  = np.linalg.norm(u - v, axis=1, keepdims=True)
    cos = dot / np.maximum(np.linalg.norm(u, axis=1, keepdims=True)
                           * np.linalg.norm(v, axis=1, keepdims=True), 1e-12)
    return np.hstack([had, dif, cos, dot, l2])

class LinkPredictor:
    C_GRID = (0.1, 1.0, 10.0)

    def fit(self, Z, pairs_train, y_train, pairs_val, y_val):
        Zn  = _unit(np.asarray(Z, dtype=float))
        Xtr = _pair_features(Zn, pairs_train)
        self.scaler = StandardScaler().fit(Xtr)
        Xtr, ytr = self.scaler.transform(Xtr), np.asarray(y_train).ravel()
        Xva = self.scaler.transform(_pair_features(Zn, pairs_val))
        nq, nc  = pairs_val.shape[:2]
        pos_val = np.argmax(np.asarray(y_val), axis=1)

        best = (-1.0, None, None)
        for C in self.C_GRID:                    # model selection on validation MRR
            clf = LogisticRegression(C=C, max_iter=2000, random_state=RANDOM_SEED)
            clf.fit(Xtr, ytr)
            s    = clf.decision_function(Xva).reshape(nq, nc)
            rank = 1 + (s > s[np.arange(nq), pos_val][:, None]).sum(1)
            mrr  = float(np.mean(1.0 / rank))
            if mrr > best[0]:
                best = (mrr, clf, C)
        self.val_mrr_, self.clf, self.C_ = best
        return self

    def score(self, Z, pairs):                   # continuous, higher = more likely
        X = self.scaler.transform(_pair_features(_unit(np.asarray(Z, float)), pairs))
        return self.clf.decision_function(X).reshape(pairs.shape[:2])
\end{lstlisting}

\begin{lstlisting}[caption={TODO 3b --- \texttt{CountryClassifier}},captionpos=b]
class CountryClassifier:
    C_GRID = (0.1, 1.0, 10.0, 100.0)

    def fit(self, Z_train, y_train, Z_val, y_val):
        Xtr = _unit(np.asarray(Z_train, dtype=float))
        self.scaler = StandardScaler().fit(Xtr)
        Xtr = self.scaler.transform(Xtr)
        Xva = self.scaler.transform(_unit(np.asarray(Z_val, dtype=float)))
        best = (-1.0, None, None)
        for C in self.C_GRID:                    # selection on validation micro-F1
            clf = LogisticRegression(C=C, max_iter=5000, random_state=RANDOM_SEED)
            clf.fit(Xtr, y_train)
            acc = float((clf.predict(Xva) == np.asarray(y_val)).mean())
            if acc > best[0]:
                best = (acc, clf, C)
        self.val_micro_f1_, self.clf, self.C_ = best
        return self

    def predict(self, Z):
        return self.clf.predict(self.scaler.transform(_unit(np.asarray(Z, float))))
\end{lstlisting}

